
Our results present an unusual scientific outcome: successful infrastructure validation alongside failed hypothesis testing due to invalid experimental conditions. This section interprets our findings, establishes guidelines for model capacity thresholds in selective prediction experiments, acknowledges honest limitations, and discusses broader implications for experimental methodology in machine learning.

\subsubsection{Interpreting the Split Outcome}

The central finding of our work is that infrastructure validation and hypothesis testing are separable concerns, and that model capacity acts as a binary gate for the latter. Our 100\% extraction rate and 8.20\% Q3 population confirm that entropy-based selective prediction is technically feasible---distributions are accessible, disagreement patterns exist, and the statistical framework is sound. This holds independent of model performance and represents a validated contribution.

However, our hypothesis---that entropy correlates with prediction correctness---remains untested, not refuted. The distinction is critical. A negative result ($\rho \approx$ 0, p > 0.05) would indicate no association exists, suggesting the hypothesis is wrong. An undefined result ($\rho$ = NaN) indicates the test cannot run, suggesting the experimental conditions are wrong. Negative results refute hypotheses; undefined results invalidate experiments.

This distinction is often missed in machine learning literature, where ''model performs poorly'' is treated as evidence against a method rather than evidence of invalid test conditions. Our MUST\_WORK gate design made the invalidity explicit, terminating the hypothesis loop and forcing us to confront the methodological failure. Without explicit validity checking, we might have incorrectly concluded that entropy does not correlate with correctness, when in fact we never tested the claim.

\subsubsection{Model Capacity as an Experimental Dependency}

Our findings suggest that selective prediction experiments have a prerequisite: models must produce non-zero variance in evaluation metrics for correlation-based statistical tests to be valid. For factual QA, this translates to a model capacity threshold. GPT-2 (117M parameters) falls below this threshold on TriviaQA, producing 0\% accuracy. Based on Roberts et al. (2020) showing that models require at least 10B parameters to achieve competitive performance on TriviaQA, we hypothesize the validity threshold lies around 7B parameters, though we tested only GPT-2 and have not empirically confirmed this across intermediate scales.

We propose the following guideline for future selective prediction research on factual QA:

\textbf{Minimum Model Capacity Threshold:} For correlation-based validation of uncertainty estimation methods on factual question-answering tasks, use models with $\geq 7B$ parameters (or verify empirically that baseline accuracy exceeds 5\% on the chosen dataset).

\textbf{Rationale:}
\begin{enumerate}
\item \textbf{Variance Requirement:} Correlation tests require variance in both variables. If accuracy is 0\% or 100\%, correctness has zero variance.
\item \textbf{Statistical Power:} Even with non-zero accuracy, very low accuracy (e.g., 1-2\%) provides minimal variance and low power. 5\% accuracy yields approximately 25 correct predictions in a 500-example sample, sufficient for preliminary correlation testing.
\item \textbf{Factual Knowledge Capacity:} Scaling laws show that factual knowledge capacity increases with parameter count. Models <1B parameters perform near-chance on knowledge-intensive tasks.
\end{enumerate}

The key principle is: \textbf{check variance empirically before assuming correlation tests are valid}. A simple pre-flight check is to compute variance(correctness) and verify it exceeds zero. This takes seconds and prevents wasted experimental runs.

\subsubsection{Honest Limitations}

\textbf{Limitation 1: Core hypothesis untested}  
Our primary research question---whether entropy-based rejection outperforms max-probability-based rejection---remains unanswered. This is the stated motivation, yet we provide no evidence for or against it. Our contribution shifts from a performance claim to a methodological insight. We identify an experimental design requirement not explicitly documented in prior work and provide infrastructure validation that future work can build upon.

\textbf{Limitation 2: Single model tested}  
We tested only GPT-2, providing one data point for the capacity threshold. Our 7B threshold is extrapolated from scaling laws literature \citep{Roberts2020}, not empirically confirmed. Our MUST\_WORK gate correctly identified the failure mode and terminated the hypothesis loop, demonstrating that the gate-based design serves its intended purpose. A multi-model capacity study is valuable future work but not essential to our methodological contribution.

\textbf{Limitation 3: Single dataset (TriviaQA only)}  
We did not test SQuAD or Natural Questions, limiting generalization claims about factual QA broadly. Infrastructure validation on one dataset is sufficient to demonstrate technical feasibility. Multi-dataset evaluation was planned for mechanism-testing sub-hypotheses, which were correctly blocked by the gate failure.

\textbf{Limitation 4: Unverified assumptions}  
Our assumptions about entropy-max-prob independence, top-5 multi-modality, and threshold coverage control remain untested. Assumption testing was planned for mechanism sub-hypotheses, which were blocked by the foundation failure. Testing assumptions when the foundation is unconfirmed would be premature.

\subsubsection{Broader Impact and Implications}

\textbf{For Selective Prediction Research:}  
Future studies must report model capacity explicitly (parameter count, baseline accuracy on chosen dataset) and verify that evaluation metrics exhibit non-zero variance before claiming negative results. Reviewers should ask: ''Is the test mathematically valid?'' before ''Did the method fail?''

\textbf{For Uncertainty Quantification:}  
Infrastructure validation (ours: extraction feasibility, disagreement patterns) can succeed independently of hypothesis validation. This allows incremental progress: even if entropy does not outperform max-probability, we know extraction is feasible and disagreement cases exist, providing a foundation for alternative approaches.

\textbf{For Experimental Methodology:}  
The distinction between undefined tests (invalid conditions) and negative tests (valid conditions, hypothesis unsupported) should be standard in machine learning reporting. A simple variance check prevents misinterpretation of methodological failures as substantive findings.
